\chapter{PENDAHULUAN}

\section{Latar Belakang}
% Ubah paragraf-paragraf berikut sesuai dengan latar belakang dari tugas akhir
Kebutuhan energi listrik rumah tangga terus meningkat seiring pertumbuhan
penduduk dan perubahan gaya hidup, sementara ketergantungan terhadap energi
fosil mendorong peralihan menuju sumber energi terbarukan yang lebih ramah
lingkungan. Salah satu solusi yang paling banyak diadopsi pada skala rumah
tangga adalah pembangkit listrik tenaga surya (PLTS) atap. Penggunaan PLTS
untuk kebutuhan rumah tangga mengalami tren peningkatan sebesar 2{,}26\% dalam
lima tahun terakhir \parencite{article}. Dalam sistem PLTS atap, inverter
memegang peran penting karena berfungsi mengubah arus searah (DC) yang
dihasilkan oleh panel surya menjadi arus bolak-balik (AC) yang dapat digunakan
oleh peralatan rumah tangga. Tanpa inverter, energi dari panel surya tidak
dapat dimanfaatkan secara langsung oleh perangkat listrik dalam rumah tangga
\parencite{10695139}.

Di antara berbagai beban listrik rumah tangga, pompa air merupakan salah satu
perangkat yang paling umum digunakan. Pompa ini berfungsi mengalirkan air dari
tandon atau sumur menuju berbagai titik pemakaian seperti kamar mandi, dapur,
dan sistem irigasi \parencite{SENTHILKUMAR2020303}. Pada instalasi air yang
melibatkan banyak titik keran, pembukaan keran secara serentak sangat
memengaruhi tekanan dan debit air yang disuplai oleh pompa. Ketika hanya sedikit
keran yang dibuka, pompa menghasilkan tekanan yang sangat tinggi, dan sebaliknya
tekanan akan turun drastis ketika banyak keran dibuka secara bersamaan
\parencite{anakub}. Ketidakstabilan tekanan ini menyebabkan distribusi air
menjadi tidak nyaman dan tidak efisien, serta dalam
jangka panjang dapat mempercepat keausan komponen pompa.

Beberapa solusi telah dikembangkan untuk menjaga tekanan pompa air tetap
konstan. Salah satunya adalah \emph{automatic pressure control} (APC) yang
bekerja dengan prinsip buka-tutup aliran berdasarkan ambang tekanan tertentu.
Selain itu, terdapat pula pendekatan yang menggunakan TRIAC \emph{driver}
berbasis kendali PID untuk mengatur daya yang masuk ke motor pompa sehingga
kecepatan dan tekanannya dapat dikontrol \parencite{anakits}.  
Namun, kendali konvensional seperti PID bersifat \emph{model-based}, yaitu
memerlukan pemodelan matematis \emph{plant} beserta penyetelan (\emph{tuning})
parameter secara manual agar diperoleh respons yang baik. Pada sistem tekanan
pompa yang nonlinier dan berubah-ubah terhadap jumlah keran yang terbuka,
pemodelan dan penyetelan tersebut memerlukan perhitugan dan model matematis.

Pada penelitian ini akan dilakukan pendekatan yang berbeda yaitu dengan kontrol (\emph{neural network}).
Dimana, pendekatan berbasis kecerdasan buatan, khususnya jaringan saraf
tiruan (\emph{neural network}), bersifat \emph{data-driven}. Alih-alih menurunkan
model matematis \emph{plant}, \emph{neural network} mempelajari pola hubungan
masukan--keluaran yang kompleks dan nonlinier langsung dari data operasi melalui
proses pelatihan (\emph{training}). Dengan demikian, kendali dapat diperoleh tanpa
memerlukan model matematis \emph{plant} secara eksplisit maupun penyetelan
parameter secara manual. Karakteristik inilah yang membuat \emph{neural network}
sesuai untuk mengendalikan sistem tekanan pompa sehingga kendali tekanan dapat
dilakukan lebih adaptif.

\section{Rumusan Masalah}
% Ubah paragraf berikut sesuai dengan rumusan masalah dari tugas akhir
Mengacu parmasalahan yang telah dipaparkan pada latar belakang, 
maka rumusan masalah yang terdapat dalam
penelitian ini adalah sebagai berikut.
\begin{enumerate}
    \item Bagaimana proses perancangan sistem kontrol tekanan pada pompa air yang terintegrasi dengan inverter?
    \item Bagaimana menerapkan \emph{neural network} sebagai sistem kendali tekanan pada pompa air?
    \end{enumerate}
    
\section{Batasan Masalah}

Batasan masalah dari penelitian ini adalah sebagai berikut.
\begin{enumerate}
    \item Desain dan simulasi dilakukan menggunakan software PSIM.
    \item Semua simulasi dan perhitungan dilakukan dalam mode ideal dan menggunakan beban dinamis. 
    \item Input dari inverter akan menggunakan baterai 12V ataupun \emph{power supply}
    \item Jenis pompa yang digunakan adalah pompa air 1 fasa
    \item Keran yang digunakan sebanyak empat buah
    \end{enumerate}

\section{Tujuan}

% Ubah paragraf berikut sesuai dengan tujuan penelitian dari tugas akhir
Penelitian ini bertujuan untuk merancang sistem inverter 
yang mampu mengendalikan tekanan pompa air secara stabil 
terhadap variasi beban dengan menggunakan algoritma \emph{neural network}.

\section{Manfaat}

% Ubah paragraf berikut sesuai dengan tujuan penelitian dari tugas akhir
Hasil yang diperoleh dari penelitian ini diharapkan dapat memberikan 
manfaat sebagai berikut:

\begin{enumerate}
    \item Dapat digunakan sebagai referensi untuk pengembangan \textit{solar water pump inverter}
    \item Dapat menjadi salah satu solusi bagi pengguna PLTS atap yang memiliki pompa air
    \end{enumerate}
    




